\section*{Lab-oppsettet}
\setcounter{section}{2}
\addcontentsline{toc}{section}{Lab-oppsettet}

\setcounter{subsection}{0}
\subsection{Lab-oppsettet}

Lab-oppsettet består av et lukket vannkretsløp der en sentrifugalpumpe driver vannet fra tanken og frem til Pelton-turbinen. Vannstrømmen reguleres med en ventil, mens dyseåpningen justeres med nålposisjonen. I forsøket ble nålposisjon 7 og 3 undersøkt.

De viktigste måleverdiene i forsøket var volumstrøm, innløpstrykk til turbinen, turbinturtall og dreiemoment på turbinakselen. Volumstrømmen ble målt med strømningsmåler, trykket ble målt ved turbinens innløp, turtallet ble registrert med turtallssensor, og dreiemomentet ble målt ved hjelp av bremsesystemet og kraftsensoren.

Belastningen på turbinen ble endret ved å stramme bremsereimen rundt bremsetrommelen. Økt belastning gir høyere dreiemoment, men lavere turtall. Dette gjør det mulig å undersøke hvordan virkningsgraden varierer med driftspunktet.

De nummererte komponentene er illustrert i figur \ref{fig:lab}.
